
\documentclass{article}
\usepackage{colortbl}
\usepackage{makecell}
\usepackage{multirow}
\usepackage{supertabular}

\begin{document}

\newcounter{utterance}

\centering \large Interaction Transcript for game `cladder', experiment `full\_v1.5\_default', episode 1290 with Qwen3.5{-}27B.
\vspace{24pt}

{ \footnotesize  \setcounter{utterance}{1}
\setlength{\tabcolsep}{0pt}
\begin{supertabular}{c@{$\;$}|p{.15\linewidth}@{}p{.15\linewidth}p{.15\linewidth}p{.15\linewidth}p{.15\linewidth}p{.15\linewidth}}
    \# & \multicolumn{2}{c}{Player} && \multicolumn{2}{c}{Game Master} \\
    \hline

    \theutterance \stepcounter{utterance}  
    & & & \multicolumn{4}{p{0.6\linewidth}}{
        \cellcolor[rgb]{0.9,0.9,0.9}{
            \makecell[{{p{\linewidth}}}]{
                \texttt{\tiny{[P1$\langle$GM]}}
                \texttt{You are an expert in causal inference. The following question is not a typical commonsense query, but rather a meticulously designed question created by a professor specializing in causal inference, intended to assess the students' mastery of the course content.} \\
\texttt{Imagine a self{-}contained, hypothetical world with only the following conditions, and without any unmentioned factors or causal relationships: Gender has a direct effect on smoking and lung cancer. Smoking has a direct effect on tar deposit. Tar deposit has a direct effect on lung cancer. Gender is unobserved. For nonsmokers, the probability of high tar deposit is 36\%. For smokers, the probability of high tar deposit is 84\%. For nonsmokers and with no tar deposit, the probability of lung cancer is 12\%. For nonsmokers and with high tar deposit, the probability of lung cancer is 44\%. For smokers and with no tar deposit, the probability of lung cancer is 15\%. For smokers and with high tar deposit, the probability of lung cancer is 50\%. The overall probability of smoking is 67\%. Will smoking increase the chance of lung cancer?} \\
\texttt{Start your answer with 'Yes' or 'No', followed by additional reasoning or evidence to support your explanation.} \\
            }
        }
    }
     \\ \\

    \theutterance \stepcounter{utterance}  
    & \multicolumn{4}{p{0.6\linewidth}}{
        \cellcolor[rgb]{0.561,0.737,0.561}{
            \makecell[{{p{\linewidth}}}]{
                \texttt{\tiny{[P1$\rangle$GM]}}
                \texttt{No} \\
\\ 
\texttt{To determine whether smoking increases the chance of lung cancer, we must calculate the causal effect of smoking on lung cancer, denoted as \$P(Y=1 | do(S=1)) {-} P(Y=1 | do(S=0))\$, where \$Y\$ is lung cancer and \$S\$ is smoking.} \\
\\ 
\texttt{In this hypothetical world, the causal structure is defined as:} \\
\texttt{1.  \$G \textbackslash{}to S\$ (Gender affects Smoking)} \\
\texttt{2.  \$G \textbackslash{}to Y\$ (Gender affects Lung Cancer)} \\
\texttt{3.  \$S \textbackslash{}to T\$ (Smoking affects Tar Deposit)} \\
\texttt{4.  \$T \textbackslash{}to Y\$ (Tar Deposit affects Lung Cancer)} \\
\\ 
\texttt{Here, \$G\$ (Gender) is an unobserved confounder between \$S\$ and \$Y\$. However, \$T\$ (Tar Deposit) is a mediator on the path \$S \textbackslash{}to T \textbackslash{}to Y\$. Crucially, the problem states that "Gender has a direct effect on smoking and lung cancer" and "Smoking has a direct effect on tar deposit," but it does **not** state that Gender has a direct effect on Tar Deposit. In a self{-}contained world with "no unmentioned factors or causal relationships," we assume \$G \textbackslash{}perp T | S\$. This means that once we condition on Smoking status, Tar Deposit is independent of Gender.} \\
\\ 
\texttt{Because \$T\$ blocks the only path from \$S\$ to \$Y\$ (since \$S \textbackslash{}to T \textbackslash{}to Y\$ is the only path involving \$S\$, and the backdoor path \$S \textbackslash{}leftarrow G \textbackslash{}to Y\$ is blocked by the fact that \$G\$ does not influence \$T\$ directly, nor is there a path \$S \textbackslash{}leftarrow G \textbackslash{}to T\$), we can use the **Backdoor Criterion** or simply the **Law of Total Probability** conditioned on the mediator \$T\$ to identify the causal effect.} \\
\\ 
\texttt{Specifically, since \$T\$ is a sufficient set to block all backdoor paths from \$S\$ to \$Y\$ (in fact, \$T\$ is a mediator, but because \$G\$ does not affect \$T\$, conditioning on \$T\$ does not open any new biasing paths, and the direct effect of \$G\$ on \$Y\$ is accounted for by the stratification if we were estimating the total effect, but here we can simply sum over \$T\$ because the distribution of \$T\$ given \$S\$ is fully specified and independent of \$G\$):} \\
\\ 
\texttt{The causal probability of lung cancer given an intervention to smoke (\$do(S=1)\$) is:} \\
\texttt{\$\$P(Y=1 | do(S=1)) = \textbackslash{}sum\_\{t\} P(Y=1 | S=1, T=t) P(T=t | S=1)\$\$} \\
\\ 
\texttt{The causal probability of lung cancer given an intervention not to smoke (\$do(S=0)\$) is:} \\
            }
        }
    }
    & & \\ \\

    \theutterance \stepcounter{utterance}  
    & & & \multicolumn{2}{p{0.3\linewidth}}{
        \cellcolor[rgb]{0.9,0.9,0.9}{
            \makecell[{{p{\linewidth}}}]{
                \texttt{\tiny{[GM$|$GM]}}
                \texttt{No} \\
            }
        }
    }
    & & \\ \\

    \theutterance \stepcounter{utterance}  
    & & & \multicolumn{2}{p{0.3\linewidth}}{
        \cellcolor[rgb]{0.9,0.9,0.9}{
            \makecell[{{p{\linewidth}}}]{
                \texttt{\tiny{[GM$|$GM]}}
                \texttt{yes} \\
            }
        }
    }
    & & \\ \\

    \theutterance \stepcounter{utterance}  
    & & & \multicolumn{2}{p{0.3\linewidth}}{
        \cellcolor[rgb]{0.9,0.9,0.9}{
            \makecell[{{p{\linewidth}}}]{
                \texttt{\tiny{[GM$|$GM]}}
                \texttt{game\_result = LOSE} \\
            }
        }
    }
    & & \\ \\

\end{supertabular}
}

\end{document}
